
\chapter{张量算法}

\section{张量缩并顺序}\label{sec:opt_contr}

在前面的章节中，我们画了\emph{大量}张量网络。虽然理论层面对某些应用已经足够，但确实存在一些研究领域，主要是量子线路模拟这类计算型领域，需要真正计算张量网络背后的张量。
这通过所谓的\emph{张量缩并}来完成。
为了理解它的含义，考虑下面这个 2-张量图，其中我们缩并两个张量 $A$ 和 $B$：
\begin{center}
    \resizebox{0.44\textwidth}{!}{
        \input{figures/contraction_1}
    }
\end{center}
注意，我们还标出了各个维度的大小，也就是说，$A$ 是一个 $2\times 2\times 4$ 张量，而 $B$ 是一个 $4 \times 2$ 矩阵。
现在，\emph{缩并成本}定义为缩并过程中执行的 FLOPS 数；按照惯例，这里指的是标量乘法次数。\footnote{这假设该操作采用朴素实现，也就是在各个维度上使用嵌套循环。}
在这个例子中，缩并成本为 $2 \times 2 \times 4 \times 2 = 32$，也就是说，只需把维度大小相乘。

前面的例子相当小。
然而，实际中的张量网络往往包含数百个张量。
自然，这些张量也需要缩并成单个张量。
有趣的地方来了：执行这些缩并的\emph{顺序}可能会对运行时间产生巨大影响。
为了对此建立直觉，考虑一个扩展的 3-张量图例子：
\begin{center}
    \resizebox{0.35\textwidth}{!}{
        \input{figures/contraction_2}
    }
\end{center}
如果先缩并 $A$ 和 $B$（再把所得张量与 $C$ 缩并），成本为 $2^2 \times 4\times 2 + 2^3 \times 1 \times 2^2 = 32 + 32 = 64$ 次乘法；而如果一开始就缩并 $B$ 和 $C$，总共会得到 $2^3 \times 1 \times 2^2 + 2^2 \times 4 \times 2^3 = 32 + 128 = 160$ 次乘法。因此，第一种顺序要好得多。

于是很自然地会问：\emph{我们总能找到最优缩并顺序吗？}

\subsection{算法}

下面总结一些用于寻找良好缩并顺序的知名结果和框架。

\subsubsection{最优算法}

事实上，为任意形状的张量网络寻找\emph{最优}缩并顺序相当困难；更准确地说，它是 NP-hard 的~\cite{ordering_np_hard}。有一个著名的精确算法运行时间为 $O(3^n)$~\cite{ordering_exact_algo}，其中 $n$ 是网络中的张量数量。它能找到相对于总缩并成本而言的最优缩并顺序。如果关心的是最小化最大中间张量的大小，也就是优化执行过程中的内存使用，那么可以在 $O(2^n n^3)$ 时间内更快地完成~\cite{dpconv}。

好消息是，对于某些受限形状的张量网络，确实存在高效算法。经典例子是矩阵链问题的动态规划解法~\cite{cormen}，它其实就是\emph{我们}的问题，只不过只针对矩阵。朴素算法运行时间为 $O(n^3)$，但可以实现为 $O(n^2)$ 时间~\cite{mc_yao}（甚至可以做到 $O(n \log n)$ 时间~\cite{mc_nlogn_1, mc_nlogn_2}）。另一类存在多项式时间算法的形状是\emph{树}张量网络~\cite{ttn_size, ttn_cost}。

另一种重要的缩并顺序优化方法，是对张量网络的\emph{线图}表示做树分解~\cite{markov_shi, dudek_tree_decomp, roman_tree_decomp}。特别地，这会得到一个缩并顺序，其最大中间张量秩等于该树分解的树宽。粗略地说，树宽衡量一个图有多像树。这并不能直接解决我们的问题，因为寻找最小树宽的树分解本身也是困难的~\cite{treewidth_approx}。两个用于寻找良好树分解的知名框架是 \texttt{QuickBB}~\cite{quick_bb} 和 \texttt{FlowCutter}~\cite{flow_cutter_1, flow_cutter_2}。

\subsubsection{尽力而为算法}

% TODO: @Thomas, refer to the `einsum` section.
然而，一旦我们想缩并\emph{任意}网络形状，能做的最好事情就是退回到启发式或近似方法。两个知名框架是 \texttt{opt\_einsum}~\cite{opt_einsum} 和 \texttt{cotengra}~\cite{cotengra}，它们旨在优化任意 einsum 的缩并顺序（也称为“缩并路径”（contraction path））：对于最优算法会过慢的张量网络，\texttt{opt\_einsum} 会应用一种特设的贪心算法，而 \texttt{cotengra} 使用超图划分~\cite{kahypar}，并结合一种后来得到改进的 Bayesian 优化方法~\cite{jena_tn_cut}。其他借鉴数据库查询优化的算法则在 \texttt{netzwerk}~\cite{ttn_cost} 中实现。另一种选择是\emph{学习}最佳缩并顺序，例如使用强化学习~\cite{meirom_tn_rl}。

自然，上述启发式方法并不带有最优性保证。存在一个运行时间为 $O^*(2^n / \eps)$ 的 $(1+\eps)$-近似算法，可以最小化中间张量大小之和~\cite{dpconv}。

值得一提的是 einsum 的 SQL 视角~\cite{einsum_as_sql}：它允许在任意数据库引擎上执行整个张量网络缩并。事实上，对于某些 einsum 类，在 Hyper~\cite{hyper} 等先进数据库引擎上运行缩并，可能比使用事实标准的 \texttt{numpy} 数组实现快得多。
